\chapter{Combinatorics 4.7}
\label{ch:250407}

\section{a combinatorial equality}

\[
\sum_{r=0}^{50}(-1)^r{101\choose 51+r}{50+r\choose 50}=1
\]

寻找一种组合方法，其代数做法在 \href{250324.md}{Lecture 06} 中.

101选51+r打勾，最后一个画圈，前面50+r个分方块三角，有 ${101\choose 51+r}{50+r\choose 50}$ 种方法.

换顺序，先选50个三角+1个圈，按顺序；再在圈前面任意选r个方块，r奇数，那么是坏的，偶数是好的。按照圈的位置分，如果圈前面选完三角还有空，那么r奇偶抵消；否则好的比坏的多1.

$\sum_{r=0}^{50}(-1)^r{101\choose 51+r}{50+r\choose 50}=\sum_{k=51}^{101}\sum_{r=0}^{k-1-50}(-1)^r{k-1\choose 50}{k-1-50\choose r}$

\section{flipping a coin}

$p\in(0,1)$ 均匀分布，求100次中50次成功的概率

传统：微积分

成功=$X(\omega)<p$, 因此这个概率等价为 rand 101 个数字，第一个数字排在第51位（or 任意位）的概率：1/101

\section{tree embedding}

\notefigure{notes/250407-01.png}

\begin{quotation}
\textbf{Thm.} $\overline{d}(G)\geq 2k-2\Rightarrow$ 任意一个 $k+1$ 阶树是 $G$ 的子图.
\end{quotation}

\begin{quotation}
\textbf{Lemma.} $\overline{d}(G)\geq 2k-2$ 那么存在 $\delta(H)\geq k$ 的诱导子图.
\end{quotation}

$m\geq n(k-1)$, 每次去掉度数 $\leq k-1$ 的点和它的邻边, $m-d(v)\geq m-(k-1)\geq(n-1)(k-1)$. 去掉的过程中，$m\geq n(k-1)$ 一直保持. 停止的时候，没有度数 $\leq k-1$ 的点，因此 $\delta(G)\geq k$.

问题：k还能再大吗？假设k=51，要求一个 $\delta\geq52$ 的诱导子图

左到右足够多的点排成一行，每个点到左边51点、右边51点连边，可以expect它平均度数大于100. 但是最左边的点度数51，去掉；因此按照这个方法会被全部去掉。

或者考虑 $K_{51,N},N\gg 51,\overline{d}=2\times\frac{51N}{51+N}\geq100$, 因此一个 $\delta\geq 52$ 的诱导子图一定不包含右边 $N$ 个点. 左边剩51个点，更不可能.

\begin{quotation}
\textbf{Lemma.} $\delta(G)\geq k$, 那么任意一个 $k+1$ 阶树是 $G$ 的子图.
\end{quotation}

容易看到，一定存在k+1阶的星星图 or $P_{k+1}$. 树任选根，每个点由于存在至少 $k$ 个邻居，总是可以逐步拓展embedding. 如果挑到x，它的朋友都在已经选的集合中，那么已经完成.

\begin{quotation}
\textbf{Erdos-Sos Conjecture.} $\overline{d}(G)\geq k\Rightarrow$ 任意一个 $k+1$ 阶树是 $G$ 的子图.
\end{quotation}

除了有限个反例，基本正确.

\section{some hard problems in graph coloring}

\notefigure{notes/250407-02.png}

前三个解答见 \href{250428.md}{Lecture 12}

$\chi\geq\omega+k$ 用 k 个 $C_5$

最后一个见 \href{https://math.mit.edu/~fox/MAT307-lecture13.pdf}{girth}，概率方法证明，用到了 Markov 不等式和 alteration

\section{一些基础不等式}

\[
\omega\leq\chi\leq n
\]

\[
\frac{n}{\alpha}\leq\chi\leq\Delta+1
\]

右边：归纳法假设30个点以下都满足 $\chi\leq\Delta+1$, 剩下一个点度数至多 29.

左边：每种颜色是一个独立集（k独立集$\iff$k染色），$n=\sum |D|\leq\sum\alpha=\chi\alpha$

问题：有没有图 $\chi\gg \frac{n}{\alpha}$？

第一反应：$K_n$ 的影子图，$n\gg \frac{2n}{n}$.

其实很容易做到，$K_n$ 再带上n个孤立点即可.

$\chi\ll\Delta+1$: star graph

$\chi=\Delta+1$: $K_n$, $C_{2k+1}$

奇妙的图 by xiaomin

\notefigure{notes/250407-03.png}

\section{定理 0407 系列}

\begin{quotation}
\textbf{Thm 040701.} $\chi(G)=k$, 则 $G$ 至少有一个点 $\deg\geq k-1$.
\end{quotation}

总有人被逼着染 k 色.

\begin{quotation}
\textbf{Thm 040702.} $\chi(G)=100$, 则 $G$ 至少有 100 个点 $\deg\geq99$.
\end{quotation}

把度数 $\geq99$ 的点提到前面，如果有 $<100$ 个这样的点，那么我可以拿着99种颜色，先给度数 $\geq 99$ 的点一人一个. 后面的点，因为度数 $\leq 98$, 用 99 个颜色一定可以染完, 得到 $\chi\leq99$, 矛盾.

归纳法：按照 $G$ 的阶数归纳，归纳起点是 $K_{100}$. 如果对任意 $|G|<n$ 成立，那么 $|G|=n,\chi(G)=100$ 时，若存在一个度数 $\deg v\leq98$, 那么 $\chi(G-v)=100$, $v$ 由于最多只看到98种颜色，它的染色不会影响 $G$ 的染色. 由归纳假设 $G-v$ 至少100个点 $\deg\geq k-1$.

\begin{quotation}
\textbf{Thm 040703'.} $\chi(G)=100$, $\phi$ 最优，则 $G$ 每种颜色都有一个点 $\deg\geq 99$.
\end{quotation}

\begin{quotation}
\textbf{Thm 040703.} $\chi(G)=100$, $\phi$ 最优，则 $G$ 每种颜色都有一个点和其他所有颜色相邻.
\end{quotation}

否则可以把这个颜色给分裂掉.

\section{贪心染色}

每次找最大独立集染色，并且去掉. 这个算法对应一种假设：

\begin{quotation}
\textbf{Conj.} $W\subseteq G$ 为最大独立集，则有一个最优染色使 $W$ 同色.
\end{quotation}

显然错误，考虑 $P_4$.

\begin{quotation}
\textbf{Conj.} 给定图 $G$, 存在最大独立集 $W\subseteq G$, 存在一个最优染色使 $W$ 同色.
\end{quotation}

$K_n$ 影子图

\section{平面直线交点形成的图}

\notefigure{notes/250407-04.png}

$\Delta=4\Rightarrow\chi\leq 5$

四色定理$\Rightarrow\chi\leq4$ (大炮打蚊子)

\begin{quotation}
\textbf{Thm. (K. Wagner, 1936; I. Fáry, 1948; Sherman K. Stein, 1951, independently)} 任意一个简单平面图总可以用直线段画出来.
\end{quotation}

或者 Brooks Thm

\begin{quotation}
\textbf{Thm. (Brooks)} Connected graph in which every vertex has at most $\Delta$ neighbors, the vertices can be colored with only $\Delta$ colors, except for two cases, complete graphs and cycle graphs of odd length, which require $\Delta+1$ colors.
\end{quotation}

这启发我们，无脑染色时的“脑子”只需要保证每个点前面有足够少的邻居即可.
